% Section 7 -- Where tone lives. Budget: 1.00 page, holds Figure 5.
% Follows docs/DESIGN_DECISIONS.md section 9 (two-condition decodability, positions, the
% legal V3 patching pair, two readouts, position groups; bi mat / bi mat is a V4 pair only).
% STATUS: hypotheses only. No model has been run. Every result is a \placeholder.
\section{Where Tone Lives}
\label{sec:tone}

Behavioural failure on \var{3} can arise because the model never represents the tone of a syllable it received in pieces, because it represents it but does not move it to the answer position, or because it moves it and cannot use it. The two Gemma~3 models that run in fp32 on a free GPU let us tell these apart.

\paragraph{Probes.}
Following \citet{2025-findings-emnlp-719}, we train layer-wise linear probes on the residual stream to decode a syllable's tone (and, secondarily, onset, coda, nucleus and rime) at four positions: the onset fragment of a split syllable, the syllable's last token, the bare tone-mark token that NFD sometimes produces, and the first token after the syllable, the one position where NFC and NFD are compared on equal terms. Stimuli are legal syllables balanced by tone, including legal non-word syllables so that tone is not confounded with lexical identity, in several carrier sentences with a tone-neutral word after the slot; train and test syllables are disjoint, with a second split held out by rime. Tone counts as \emph{linearly decodable} at a layer and position only when two conditions hold: \emph{selectivity}, accuracy at least $\delta_{\mathrm{sel}}$ above a control probe trained on random labels per syllable type \citep{D19-1275}, with $\delta_{\mathrm{sel}} = 0.15$ provisional and frozen from the pilot; and \emph{excess} over a structural baseline, the same probe fitted on one-hot token ids plus the coda class, because stop codas admit only two tones and, under NFC, token count itself predicts tone. Probes show that information is present, not that it is used \citep{2022-cl-1-7}; we report a layer-0 baseline curve, and, on syllables that are single tokens, the layer at which tone accuracy first rises to half its maximum gain over the embedding, the tone counterpart of the reconstruction the character-level work reports.

\paragraph{Patching and steering.}
Activation patching locates where transplanted states change the output \citep{meng-etal-2022-locating,heimersheim-nanda-2024-use}. A pair is two prompts identical token for token except one syllable whose tone differs, both syllables attested and of equal segmental make-up; the varying syllable is the \emph{second} input syllable, because under \var{3} the first output syllable takes its tone and the two gold answers then diverge at their first token (clean \vi{bí mà} $\rightarrow$ \vi{bì má}, corrupt \vi{bí mạ} $\rightarrow$ \vi{bị má}; readout \vi{bì} vs.\ \vi{bị}, single pieces; pseudo-phrases, \placeholder{[NATIVE-CHECK]}). Pairs are kept only if their tokenizations align, the readout pieces differ in the tone mark alone, and the model's clean answer is correct with a logit gap of at least one nat; a model that loses every pair at the last filter has no manipulation readout, which is itself a result. We cache the clean residual stream, run the corrupt input with layer $\ell$ overwritten at a position group, and read the logit difference between the two answers at the first diverging token, normalized by the clean--corrupt gap. Two readouts separate perception from manipulation: a \emph{perception} readout asks which tone the syllable carries and reads a digit; the \emph{manipulation} readout is the \var{3} prompt itself. Position groups are the varying syllable and each of its tokens separately, the fixed syllable, the instruction, the answer marker and the final position; patching everything from the varying syllable onward must recover fully. The $2 \times 2$ of readout by source position distinguishes intact perception, intact manipulation, a copy-type failure (the answer tracks the first input syllable), and an ignore-type failure, in which the perception readout recovers where the manipulation readout does not. A difference-in-means direction between two tone classes, the baseline of CausalGym \citep{2024-acl-long-785}, tests steering against a random direction of equal norm and an onset-class direction that must not flip tone. Gemma Scope~2 sparse autoencoders are optional and must beat the probe baseline to be reported \citep{kantamneni25a}.

\paragraph{Pre-registered hypothesis.}
\hyp{5} (\emph{perception without manipulation}): in a Gemma~3 model whose \var{3} accuracy is below 25\%, tone meets both decodability conditions at a non-local position, the token after the syllable, at some layer in the first half of the network; and patching shows the perception readout recovering where the manipulation readout does not. Such a dissociation places the failure in manipulation rather than perception. Under NFD, tone is expected to become decodable at an earlier layer than under NFC because the tone mark is its own token. A clean localization without the dissociation would be equally informative.

\paragraph{Results.}
\placeholder{Figure~\ref{fig:probes}: layer-wise selectivity and excess for tone at the four positions under NFC and NFD in Gemma~3 1B and 4B, with the layer-0 baseline; the patching heatmap (layer $\times$ position group) for both readouts with recovery. Text: the first decodable layer per model, position and encoding; the patching bands for perception and manipulation; steering flip rates against the two null directions; pair retention at every filter; verdict on \hyp{5}.}

\begin{figure}[t]
  \centering
  \fbox{\parbox{0.95\columnwidth}{\centering\vspace{2.5cm}\placeholder{Figure 5: probe selectivity and excess by layer (NFC vs.\ NFD, four positions) and patching heatmap for the two readouts, Gemma~3 1B and 4B}\vspace{2.5cm}}}
  \caption{\placeholder{Left: probe selectivity (accuracy minus control) and excess over the structural baseline for tone at each layer of Gemma~3 1B and 4B, under NFC and NFD, at the syllable's last token and at the token after it. Right: recovery of the clean logit difference when the residual stream at layer $\ell$ and a position group is patched from the clean run, for the perception and the manipulation readout.}}
  \label{fig:probes}
\end{figure}
